% Pursuing Stacks (cotes 134-2 to 134-8): archive cote 134-2, site page 45 (PDF page 46).
% appendix (letters to Breen). Draft compared: ps1.tex:1214-1221.
% AI-transcribed by Claude Opus 5.5 (claude-opus-5-5, adaptive thinking, effort high), Message Batches API, batch ps1 (msgbatch_01Rv5hCaNTtUt4JHx1yqoHPL), 2026-11-12; not yet checked by eye.
\scanpage{134-2/45}
\mspage{\hand{30}}

en homotopie. Ainsi, l'espace des lacets $\hand{\Omega}$(K.,x) correspond manifestement
à la (n-1)-catégorie \uline{Hom}(x,x) formée des i-flèches de C (1 $\hand{\leq}$i$\hand{\leq}$ n\xout{\unsure{-1}})
dont la 0-origine et la 0-extrém\xout{t}ité sont x, réindexées en les appelant
(i-1)-flèches. Je \xout{\illegible{1}} n'aperçois pas à vue de nez un joli candidat pour la
suspension en termes de n-catégories. Par contre le Hom.(K.,K.') doit cor-%
respondre au \uline{Hom}(C,C'), qui est une n-catégorie quand C,C' en sont. La
``fibre homotopique'' d'une application semi-simpliciale f: K.$\hand{\to}$K.'(trans
formée d'abord, pour les besoins de la cause, en une fibration de Serre
par le procédé bien connu de Serre-Cartan) correspond sans doute à l'opéra-%
tion bien familière de produit \hand{(}n\inserted{+1)}-fibré (du moins les cas n=0,1 sont bien
familiers !) \inserted[typed]{\xout{(n+1)}} $C\times_{C'}C''$ pour des n-foncteurs C $\hand{\to}$C' et C''$\hand{\to}$C', dans le cas
où C'' est la n-catégorie ponctuelle\hand{,} donc la donnée de C''$\hand{\to}$C' correspond à
la donnée d'un 0-objet de C'. Les espaces K($\hand{\pi}$,n) ont une interprétation
évidente comme n-gerbes liées par $\hand{\pi}$. Enfin, on voit aussi poindre l'analo-%
gue du dévissage de Postnikoff d'un ensemble semi-simplicial - mais la
façon dont je l'entrevois (vue ma prédilection pour les topos) passe par
la notion de topos classifiant d'un\xout{\unsure{e}} n-groupoïde (généralisant de façon
évidente le topos classifiant d'un groupe). En termes de cette notion,
on peut, il me semble, interpréter un n-groupoïde en termes d'un (n-1)-grou
poïde (savoir son tronqué), \uline{muni} d'une n-gerbe sur le topos classifiant,
liée par $\hand{\pi}_n$ \hand{(``\unsure{fordu|tordu}'' bien sûr par l'action du $\pi_1$ ...)}\hand{.}

\margin{espace}
\hand{(App.\,4)} Bien sûr, il faut relativiser encore tout le yoga qu'on vient de décri-%
re\hand{,} au dessus d'un topos quelconque X. Il s'agirait donc de mettre en rela-%
tion et d'identifier, dans un\xout{\unsure{e}} certaine mesure, d'une part l'algèbre homo-%
topique sur X construite en termes de la notion de faisceau semi-simplicial
sur X , d'autre part l'algèbre catégorique sur X construite en termes de
la notion de n-champ en groupoïdes (n $\hand{\geq}$ 0 fini ou infini). On espère que
la notion d'image inverse de faisceau semi-simplicial par un morphisme de
topos f: X$\hand{\to}$X' (qui est évidente) \xout{\illegible{1}} correspond à la notion \over{\unsure{p}}{p}lus
subtile d'image inverse de n-champs; et inversement, \xout{\illegible{1}} la notion évidente
d'image directe de n-champs par f devrait correspondre à une notion plus
subtile d'image directe\inserted{$Lf_*(K.)$,} d'un faisceau semi-simplicial, construit sans doute
dans l'esprit des foncteurs dérivés à partir de la notion naive (mais on
hésite s'il faut mettre Lf$\hand{_*}$ ou Rf$\hand{_*}$)... Les dévissages à la P\over{\unsure{q}}{o}stnikoff
\xout{\unsure{en termes de n-champs}} doivent avoir encore une interprétation remarquable-%
ment simple en termes de n-champs. Comparer à la remarque de Giraud qu'un
(1)-champ en groupoides sur X peut s'identifier au couple d'un faisceau
$\hand{\pi}_0$ sur X, et d'une (1)-gerbe sur le topos induit X$_{\hand{/\pi}_0}$ (dont le lien,
comme de juste, devrait être noté $\hand{\pi}_1$ !). D'ailleurs, dans le cas des
1-champs en groupoides, la traduction de ces animaux en termes de topos

\margin{45}
